% Section 8.6, from the summary table of lectures/L08/appendix/a_theory.md and from
% lectures/L08/README.md: the objectives, the evaluation questions and the next lesson.

\section{Sammanfattning}
\label{c8:sec:sammanfattning}

\begin{booktable}{@{}p{12.5em}L@{}}
  \thead{Begrepp} & \thead{Förklaring} \\
  \midrule
  Funktion & Varje $x$ ger exakt ett $f(x)$ \\
  Definitionsmängd & Tillåtna $x$-värden \\
  Värdemängd & Möjliga $f(x)$-värden \\
  Linjär funktion & $f(x) = kx + m$ \\
  Exponentialfunktion & $f(x) = C \cdot a^x$ \\
\end{booktable}

\needspace{6\baselineskip}
\subsection*{Det här ska du kunna}
Efter det här kapitlet ska du:
\begin{itemize}
  \item Ha repeterat de viktiga begreppen om funktioner från grundskolan och Matematik~1.
  \item Förstå begreppet funktion.
  \item Känna igen de vanliga funktionstyperna.
  \item Kunna använda funktioner och grafer för att lösa problem.
\end{itemize}

\testadigsjalv
\begin{enumerate}
  \item Ge ett exempel på en funktion med begränsad definitionsmängd och förklara varför den är
    begränsad.
  \item Beräkna $f(2)$ och $f(-1)$ för $f(x) = 2x^2 - x + 3$.
\end{enumerate}

\needspace{6\baselineskip}
\subsection*{Nästa kapitel}
\Kapref{9} handlar om trigonometriska funktioner: amplitud, period, fas och vinkelmått.
